\section{Direct Sums and Free Modules}

\begin{definition}
    Let $(M_i)_{i \in I}$ be a family of $R$-modules. The \emph{direct product} of $(M_i)_{i \in I}$ is the module
    \[
        \prod_{i \in I} M_i
        =
        \left\{ \varphi\colon I \to \bigcup_{i \in I} M_i \,\middle|\, \varphi(i) \in M_i \text{ for all } i \in I \right\},
    \]
    with componentwise addition and scalar multiplication:
    \[
        R \times \prod_{i \in I} M_i \to \prod_{i \in I} M_i, \quad (r, \varphi) \mapsto (r \varphi(i))_{i \in I}
    \]
    The \emph{direct sum} of $(M_i)_{i \in I}$ is the submodule
    \[
        \bigoplus_{i \in I} M_i
        =
        \left\{ \varphi \in \prod_{i \in I} M_i \,\middle|\, \#\{\, i \in I \mid \varphi(i) \neq 0 \,\} < \infty \right\}.
    \]
    These two modules agree when $|I| < \infty$. In the category of $R$-modules, the direct product is the categorical product and the direct sum is the categorical coproduct.
\end{definition}

\begin{theorem}
    Suppose that $M_1, \ldots, M_n$ are $R$-submodules of $M$ such that $M = M_1 + \cdots + M_n$. Let $I = [n]$. Then the following conditions are equivalent:
    \begin{enumerate}
        \item The map $\varphi\colon \bigoplus_{i \in I} M_i \to M$ given by $\varphi((m_i)_{i \in I}) = \sum_{i \in I} m_i$ is an isomorphism.
        \item Any element of $M$ can be uniquely written as a sum of elements of $M_i$.
        \item The zero element can be uniquely written as a sum of elements of the $M_i$.
        \item For each $i \in I$, one has
              \[
                  M_i \cap \sum_{j \neq i} M_j = 0.
              \]
    \end{enumerate}
\end{theorem}

\begin{proof}
    We prove
    \[
        (1)\Rightarrow(2)\Rightarrow(3)\Rightarrow(4)\Rightarrow(1).
    \]

    \noindent\textbf{$(1)\Rightarrow(2)$.}
    If $\varphi$ is an isomorphism, then for every $m\in M$ there exists a unique
    \[
        (m_1,\ldots,m_n)\in \bigoplus_{i=1}^n M_i
    \]
    such that
    \[
        m=\varphi(m_1,\ldots,m_n)=m_1+\cdots+m_n.
    \]

    Hence every element of $M$ can be written uniquely as a sum of elements of the $M_i$.

    \noindent\textbf{$(2)\Rightarrow(3)$.}
    Apply (2) to the element $0 \in M$.

    \noindent\textbf{$(3)\Rightarrow(4)$.}
    Fix $i\in I$, and let
    \[
        x\in M_i\cap \sum_{j\ne i} M_j.
    \]
    Then $x \in M_i$, and there exist $m_j \in M_j$ for $j \ne i$ such that
    \[
        x = \sum_{j\ne i} m_j.
    \]
    Therefore
    \[
        0 = x - \sum_{j\ne i} m_j
    \]
    is a decomposition of $0$ as a sum of elements from the submodules $M_1,\ldots,M_n$.
    By the uniqueness in $(3)$, this decomposition must be the trivial one. In particular,
    \[
        x = 0.
    \]
    Hence
    \[
        M_i\cap \sum_{j\ne i} M_j = 0.
    \]

    \noindent\textbf{$(4)\Rightarrow(1)$.}
    Since
    \[
        M = M_1+\cdots+M_n,
    \]
    the map $\varphi$ is surjective. It remains to show that $\varphi$ is injective.

    Suppose that
    \[
        \varphi(m_1,\ldots,m_n) = 0,
    \]
    that is,
    \[
        m_1+\cdots+m_n = 0
        \qquad \text{with } m_i \in M_i.
    \]
    For each $i \in I$, we have
    \[
        m_i = -\sum_{j\ne i} m_j \in \sum_{j\ne i} M_j.
    \]
    Since also $m_i \in M_i$, it follows that
    \[
        m_i \in M_i\cap \sum_{j\ne i} M_j = 0.
    \]
    Thus $m_i = 0$ for all $i$, so
    \[
        \ker\varphi = 0.
    \]
    Therefore $\varphi$ is injective, and hence an isomorphism.

    This proves that the four conditions are equivalent.
\end{proof}

\begin{definition}[free module]
    Let $M$ be an $R$-module.
    \begin{enumerate}
        \item $x_{1}, \dots, x_{n} \in M$ are linearly independent if
              \[
                  a_{1} x_{1} + a_{2} x_{2} + \cdots + a_{n} x_{n} = 0 \to a_{1} = a_{2} = \cdots = a_{n} = 0.
              \]
        \item $M$ is called free if it is generated by a linearly independent subset. This subset is called a basis of $M$.
        \item If $M$ is an $R$-module with a basis $x_{1}, \ldots, x_{n}$, then $M = \langle x_{1} \rangle \oplus \cdots \oplus \langle x_{n} \rangle$.
    \end{enumerate}
\end{definition}

\begin{example} $R \oplus \cdots \oplus R = R^{n}$.
\end{example}

\begin{theorem}
    Let $R$ be a commutative ring with $1$. Let $M$ be a finitely generated free $R$-module. Then any basis of $M$ has the same cardinality. The number of elements in a basis is called the rank of the free module.
\end{theorem}

\begin{proof}
    Let $B=\{x_i\}_{i\in I}$ be a basis of $M$.
    We first show that any basis is finite.
    Since $M$ is finitely generated, choose generators $y_1,\dots,y_n$ of $M$.
    Each $y_k$ is a finite $R$--linear combination of elements of $B$,
    so there exists a finite subset $I_0\subset I$ such that
    \[
        y_1,\dots,y_n \in \sum_{i\in I_0}Rx_i.
    \]
    Hence
    \[
        M = \sum_{k=1}^n Ry_k \subseteq \sum_{i\in I_0}Rx_i \subseteq M,
    \]
    so $M=\sum_{i\in I_0}Rx_i$.
    If $I\setminus I_0\neq\varnothing$, pick $j\in I\setminus I_0$. Then
    \[
        x_j \in M = \sum_{i\in I_0}Rx_i,
    \]
    contradicting that $B$ is linearly independent. Thus $I=I_0$ is finite.

    Now let $x_1,\dots,x_r$ and $y_1,\dots,y_s$ be two bases of $M$.
    Let $J\subset R$ be a maximal ideal and set $F:=R/J$, a field.
    Then $M/JM$ is naturally an $F$-vector space.
    Write $\overline{m}$ for the class of $m\in M$ in $M/JM$,
    and $\overline{a}$ for the class of $a\in R$ in $F$.
    We claim that $\overline{x_1},\dots,\overline{x_r}$ form an $F$--basis of $M/JM$.

    \emph{Spanning.}
    For any $m\in M$, write $m=\sum_{i=1}^r a_i x_i$. Then
    \[
        \overline{m} = \sum_{i=1}^r \overline{a_i}\,\overline{x_i},
    \]
    so $\overline{x_1},\dots,\overline{x_r}$ span $M/JM$.

    \emph{Linear independence.}
    Suppose
    \[
        \sum_{i=1}^r \overline{a_i}\,\overline{x_i} = 0 \quad \text{in } M/JM.
    \]
    Then $\sum_{i=1}^r a_i x_i \in JM$, so there exist $a'_1,\dots,a'_r\in J$ such that
    \[
        \sum_{i=1}^r a_i x_i = \sum_{i=1}^r a'_i x_i.
    \]
    Hence
    \[
        \sum_{i=1}^r (a_i - a'_i)\,x_i = 0.
    \]
    By linear independence of $x_1,\dots,x_r$, we get $a_i - a'_i = 0$ for all $i$,
    so $a_i\in J$, thus $\overline{a_i}=0$ for all $i$.
    Therefore $\overline{x_1},\dots,\overline{x_r}$ are $F$-linearly independent,
    and form an $F$--basis of $M/JM$.

    By the same argument,$\overline{y_1},\dots,\overline{y_s}$ also form an $F$--basis of $M/JM$. Hence
    \[
        r = \dim_F(M/JM) = s.
    \]
    Therefore any two bases of $M$ have the same cardinality.
\end{proof}

